\section{The guardrail under study}\label{sec:system}

\paragraph{Pattern tiers (baseline)} SafeStream's pattern tiers (regular expressions with checksum and format validators, label-anchored identifiers, a contextual scorer and an entropy tier) and its streaming engine, whose overlap rule makes chunked output equal to whole-document output for detectors that respect their declared bounds, are frozen, as T0, at the commit that opened this study; every T0 number below is only the baseline of a paired difference.

\paragraph{The gated neural tier} An optional tier runs a token-classification model on the user's machine through ONNX Runtime \cite{onnxruntime}: DistilBERT \cite{sanh2019distilbert} fine-tuned on CoNLL-2003 \cite{tjongkimsang2003conll} for person, organisation and location names, with weights dynamically quantised to INT8 \cite{jacob2018quantization}. To keep the cost of a neural model off machine text, a regular-expression \emph{gate} chooses what the model reads. A candidate is a run of two to four capitalised words, a single capitalised word after a lower-case word or a comma, or up to three words after a person cue (an honorific, ``signed by'', or a label such as ``patient:''). Each candidate is widened by 32 characters within its line, overlapping windows are merged, words after a cue are re-capitalised for the cased model, and only the merged windows reach the model. The tier declares a maximum span and left and right context of 256 characters each, which the engine adds to its overlap requirement. An optional batch mode runs up to 32 windows per call; as INT8 activation ranges are computed per call, a window's scores then depend on its batch, which depends on the chunking.

\paragraph{Name-list tier} A third optional tier, the gazetteer, needs no model: it matches capitalised words, and words after a person cue, against lists of given and family names. It is evaluated here only for egress (Section~\ref{sec:egress}).

\paragraph{Tool surface} Table~\ref{tab:contract} lists what each interface returns. The MCP server \cite{mcpspec} speaks JSON-RPC over standard input and output, and the REST server listens on loopback only. Both repeat a reply's structured content as a JSON text block, as MCP clients expect.

\begin{table}[t]
\caption{What each interface returns. Only \texttt{redact} returns input text, by design: the input with every detected value replaced.}\label{tab:contract}
\centering\footnotesize
\setlength{\tabcolsep}{3pt}\begin{tabular}{@{}l p{0.6\columnwidth}@{}}
\toprule
Interface & Fields returned \\
\midrule
\multicolumn{2}{@{}l}{\emph{MCP tools}} \\
\texttt{scan} & per finding: type, start, end, length, confidence, severity, detector \\
\texttt{explain\_detection} & as \texttt{scan}, plus a context note and detector \texttt{notes} \\
\texttt{detectors} & tier names and declared bounds \\
\texttt{health} & version, tiers, settings \\
\texttt{redact} & redacted text; counts by type \\
\midrule
REST \texttt{/scan} & per finding: start, end, type, confidence, detector \\
SARIF 2.1.0 \cite{sarif210} & rule per type; per finding: level, region (offsets, lines), a fixed message \\
\bottomrule
\end{tabular}
\end{table}

\paragraph{Redaction modes} A detected value is replaced by a fixed marker (\emph{replace}), masked except for at most its last four characters and never more than half of it (\emph{mask}), a salted digest (\emph{hash}), an HMAC-SHA-256 pseudonym \cite{rfc2104} (\emph{pseudonymize}), a deterministic AES-256-GCM token that a key holder can reverse and that fails closed on any modification (\emph{encrypt}), or a synthetic value of the same type from reserved ranges (\emph{fictionalize}). Each is a function of the value alone, which is what keeps chunked output equal to whole-document output.
